\documentclass[11pt]{article}
\usepackage{fontspec}
\usepackage[french,english]{babel}
\usepackage{amsmath}
\usepackage{amssymb}
\usepackage{amsthm}
\usepackage{mathrsfs}
\usepackage{tikz-cd}
\usepackage{enumitem}

\newcommand{\printedpage}[1]{\marginpar{\footnotesize p.~#1}}
\newcommand{\msp}[1]{\marginpar{\footnotesize [#1]}}
\newcommand{\uGamma}{\underline{\Gamma}}
\newcommand{\Ind}{\operatorname{Ind}}
\newcommand{\Ob}{\operatorname{Ob}}
\newcommand{\card}{\operatorname{card}}
\newcommand{\trois}[1]{\mathrel{\overset{#1}{\vcenter{\offinterlineskip\hbox{$\longrightarrow$}\vskip0.5pt\hbox{$\longrightarrow$}\vskip0.5pt\hbox{$\longrightarrow$}}}}}

\begin{document}

\printedpage{89}
% head: A. Grothendieck : Dérivateurs, Ch. XVIII, Édition des Universités de Montpellier II et Paris VII
% foot: Version 17 décembre 2025

The filtering case is already known (cf. II \dots). For the ``large relative to $\pi$'' case, we are going to use the

\textbf{Corollary 4.9.6.9.2.} \textit{Let $J$ be a category. For $J$ to be large relative to $\pi$, it}
\msp{page 238}
\textit{is necessary and sufficient that it be $0$-connected, and satisfy PF$_\pi$ 2 (cf. page 81), as well as the following variant of PF$_\pi$ 1$'$.}
\begin{itemize}[leftmargin=4em]
\item[\textit{PF$_\pi$ 1)}] \textit{For every object $x_0$ of $J$, $x_0\backslash J$ satisfies PF$_\pi$ 1$'$, i.e. for every family of arrows $x_0 \xrightarrow{u_\alpha} y_\alpha$ ($\alpha \in A$, $\card A \leqslant \pi$), there exists an arrow $x_0 \xrightarrow{v} z$ and arrows $v_\alpha : y_\alpha \longrightarrow z$, with $v_\alpha u_\alpha = v \ \forall\, \alpha \in A$.}
\end{itemize}

\textsc{Proof}. If $J$ is large relative to $\pi$, then PF$_\pi$ 1) follows at once from the conjunction of PF$_\pi$ 1$'$ and PF$_\pi$ 2, hence the ``only if''. Conversely, suppose that $J$ satisfies PF$_\pi$ 1, PF$_\pi$ 2, and is 0-connected, we must prove that it satisfies PF$_\pi$ 1$'$. But we already know that the conditions in question, which imply PF 1, PF 2, 0-connected, imply that $J$ is filtering. Thus if we have a family $(x_\alpha)_{\alpha\in A}$, $\card A \leqslant \pi$, choosing an $x_0 \in \Ob J$, we can for every $\alpha$ find $y_\alpha$ dominating $x_0$ and $x_\alpha$, hence arrows
\begin{tikzcd}[row sep=0.4em, column sep=1em, baseline=(\tikzcdmatrixname-2-2.base)]
x_\alpha \arrow[dr] & \\
& y_\alpha \\
x_0 \arrow[ur] &
\end{tikzcd}, we can apply PF$_\pi$ 1 to them, hence arrows $y_\alpha \longrightarrow z$, and by composing, arrows $x_\alpha \longrightarrow z$, q.e.d. % [sic] leurs appliquer -> leur appliquer, grammar

Let us now prove the lemma. It suffices (by the preceding corollary) to prove that $J$ satisfies PF$_\pi$ 1 and PF$_\pi$ 2. Now PF$_\pi$ 1 is immediate, since $\forall\, x_0$ in $J$, $x_0\backslash J \xrightarrow{\ \sim\ } y_0\backslash I$, where $y_0$ is
\msp{page 239}
the image of $x_0$ in $I$. As for PF$_\pi$ 2, if we have $x_0 \trois{u_\alpha} x_1$ in $J$, hence $g(x_0) \trois{g(u_\alpha)} g(x_1)$, we equalize in $I$ by $g(x_1) \xrightarrow{w} y$, but since $g$ is coétale, $\exists\, x_1 \xrightarrow{v} x_2$ in $J$ such that $g(v) = w$, so we have $g(vu_\alpha) = g(vu_\beta)$ for $\alpha, \beta \in A$, and since $g$ is faithful, we conclude from this that $vu_\alpha = vu_\beta$, so $v$ equalizes $(u_\alpha)_{\alpha\in A}$, q.e.d.

Let us prove b$_\pi$) $\Longrightarrow$ a$_\pi$): If $I'$ and the $i'\backslash I$ are large relative to $\pi$, then $I$ is too (and $f$ is cofinal, the $i'\backslash I$ being 0-connected, since filtering). Since we already know that $I$ is 0-connected ($I'$ being so and $f$ cofinal), it remains, by 4.9.6.9.2, to verify PF$_\pi$ 1 and PF$_\pi$ 2.

\msp{page 240}
The proof given in 1° for the similar filtering case (implication b$'$ $\Longrightarrow$ a, pages 231-232), applies mutatis mutandis. (Property F 3 has already been established in part 1°), as a consequence of b, and can therefore be used freely.)

Thus we have proved a$_\pi$ $\Longleftrightarrow$ b$_\pi$. Let us prove a$_\pi$ $\Longrightarrow$ c$_\pi$. We already know that a$_\pi$ implies that $I'$ is large relative to $\pi$, and of course condition F 1, it remains to establish the validity of condition F 3$_\pi$. By F 3, $\forall\, \alpha \in A$ we can find $j_\alpha$ in $I$ and $w^0_\alpha : i_\alpha \longrightarrow j_\alpha$, $v^0_\alpha : i \longrightarrow j$ such that % [sic] v^0_\alpha : i \longrightarrow j -> v^0_\alpha : i \longrightarrow j_\alpha, the following diagram has v^0_\alpha : i \to j_\alpha
\[
f(v_\alpha)u'_\alpha \ = \ f(w_\alpha) \ .
\]

\printedpage{90}
% head: A. Grothendieck : Dérivateurs, Ch. XVIII, Édition des Universités de Montpellier II et Paris VII
% foot: Version 17 décembre 2025

\[
\begin{tikzcd}[column sep=large, row sep=huge]
i_\alpha \arrow[rr, bend left=20, "w^0_\alpha"] \arrow[d, dash] & i \arrow[r, "v^0_\alpha"] \arrow[d, dash] & j_\alpha \arrow[r, "\lambda_\alpha"] \arrow[d, dash] & j \arrow[d, dash] \\
f(i_\alpha) \arrow[r, "u'_\alpha"] & f(i) \arrow[r, "f(v_\alpha)"] & f(j_\alpha) \arrow[r] & f(j)
\end{tikzcd}
\]

Since $I$ satisfies PF 1$_\pi$ (cf. 4.9.6.9.2), there exist $h$ in $I$ and $\lambda_\alpha : j_\alpha \longrightarrow j$, such that $\lambda_\alpha v^0_\alpha = \lambda$ is independent of $\alpha$, then we have for every $\alpha \in A$ % [sic] h -> j, the arrows \lambda_\alpha go to j
\[
f(\lambda)u'_\alpha \ \Bigl( \ = \ f(\lambda_\alpha)\underbrace{f(v_\alpha)u'_\alpha}_{f(w_\alpha)} \ \Bigr) \ = \ f(\underbrace{\lambda_\alpha w^0_\alpha}_{w_\alpha}) \ ,
\]
so $\lambda : i \longrightarrow j$ and the $w_\alpha = \lambda_\alpha w^0_\alpha : i_\alpha \longrightarrow j$ satisfy the desired conditions.

c$_\pi$ $\Longrightarrow$ b$''_\pi$. We must prove that if c$_\pi$ is satisfied, then the $i'\backslash I$ satisfy PF 1$'_\pi$.

\msp{page 241}
Let then $\xi_\alpha = (i_\alpha, u'_\alpha)$ ($\alpha \in A$, $\card A \leqslant \pi$) be a family of elements of $i'\backslash I$, since $I'$ satisfies PF 1$_\pi$ (cf. 4.9.6.9.2), there exist $j'$ in $I'$ and arrows $u''_\alpha : f(i_\alpha) \longrightarrow j'$, such that $u''_\alpha u'_\alpha = u'$ independent of $\alpha$. Since $f$ satisfies F 1, we can assume that $j'$ is of the form $f(j)$. Applying F 3$_\pi$ to the $i_\alpha$, $j$, $u''_\alpha$, we find $v : j \longrightarrow k$ and $w_\alpha : i_\alpha \longrightarrow k$, such that
\[
f(v)u''_\alpha \ = \ f(w_\alpha) \qquad \forall\, \alpha \in A \ ,
\]
hence
\[
f(w_\alpha)u'_\alpha \ = \ f(v)\,\underbrace{u''_\alpha u'_\alpha}_{u'} \ = \ f(v)u' \qquad \text{independent of $\alpha$,}
\]
say $w'$, so the $w_\alpha$ define an arrow from the $\xi_\alpha$ to $\xi = (k, w')$, q.e.d.

b$''_\pi$ $\Longrightarrow$ b$_\pi$ if $f$ is faithful. We are reduced to proving that if b$''_\pi$ is satisfied, $f$ faithful, then the categories $i'\backslash I$, which we already know satisfy PF 1$_\pi$, also satisfy PF 2$_\pi$. Let then $\xi = (i, u')$, $\eta = (j, v')$ be two objects of $i'\backslash I$, and $(u_\alpha)_{\alpha\in A}$ a family of arrows $\xi \longrightarrow \eta$ (with $\card A \leqslant \pi$), hence arrows in $I$ such that $f(u_\alpha)u' = v' \ \forall\, \alpha \in A$.
\[
\begin{tikzcd}[column sep=large]
i \arrow[r, shift left=2, "u_\alpha"] \arrow[r] \arrow[r, shift right=2] & j \arrow[r, dotted, "v"] & k
\end{tikzcd}
\]
\[
\begin{tikzcd}[column sep=large]
i' \arrow[r, "u'"] \arrow[rr, bend right=15, "v'"'] & f(i) \arrow[r, shift left=2] \arrow[r] \arrow[r, shift right=2] & f(j) \arrow[r] & f(k)
\end{tikzcd}
\]

Since $I'$ is large relative to $\pi$, we can equalize by $f(j) \longrightarrow k'$, and since $k'\backslash I \neq \emptyset$, we can assume $k'$ of the form $f(k)$. Now
\msp{page 242}
by the implication b$''$ $\Longrightarrow$ a $\Longrightarrow$ c of part 1°) (cf. page 235) and since b$''_\pi$ $\Longrightarrow$ b$''$, we already know that $f$ satisfies F 3. Applying it to $j$, $k$ and $f(j) \longrightarrow f(k)$, we can
\printedpage{91}
% head: A. Grothendieck : Dérivateurs, Ch. XVIII, Édition des Universités de Montpellier II et Paris VII
% foot: Version 17 décembre 2025
assume that $f(j) \longrightarrow f(k)$ comes from $v : j \longrightarrow k$. We shall therefore have $f(v)f(u_\alpha) = f(v)f(u_\beta)$ $\forall\, \alpha, \beta \in A$, so $f(vu_\alpha) = f(vu_\beta)$. Since $f$ is faithful, this implies that $vu_\alpha = vu_\beta$, i.e. $v$ equalizes the $u_\alpha$. Besides, $v$ can be regarded as an arrow \textit{in} $i'\backslash I$, from $(j, v')$ to $(k, f(v)v')$. Q.e.d.

\textbf{4.9.6.10.} When $I$ is \textit{ordered}, then the four conditions a$_\pi$, b$_\pi$, c$_\pi$, $b''_\pi$ are equivalent, since the $i'\backslash I$ are ordered, and it is clear that then b$''_\pi$ $\Longleftrightarrow$ b$_\pi$. % [sic] claire -> clair, grammar

\msp{page 243}

\section*{4.10\quad The indicial representation theorem: proof}

We are first going to reformulate theorem 4.9.3 in a more general form, in terms of \textit{cofibered} categories. Let
\[
\mathcal{D} \longrightarrow J \tag*{(4.10.1)}
\]
be a cofibrant functor, where $J$ is equivalent to a small category, and likewise for the fiber categories $\mathcal{D}_j$ ($j \in \Ob J$). We give ourselves an infinite cardinal $\pi$. Let
\[
\mathcal{M}_j \ = \ \Ind_\pi(\mathcal{D}_j) \qquad (j \in \Ob J) \ , \tag*{(4.10.2)}
\]
then the $\mathcal{M}_j$ are the fibers of a cofibered category $\mathcal{N}$ over $J$, which I also denote by $\Ind_\pi(\mathcal{D}/J)$, so we put
\[
\mathcal{N} \ = \ \Ind_\pi(\mathcal{D}/J) \ . \tag*{(4.10.3)}
\]
We thus have a fully faithful inclusion of cofibered categories over $J$
\[
\mathcal{D} \hookrightarrow \mathcal{N} \ ,
\]
hence an inclusion
\[
\uGamma(\mathcal{D}/J) \longrightarrow \uGamma(\mathcal{N}/J) \ . \tag*{(4.10.4)}
\]

\msp{page 244}
In the case where $\mathcal{D}$ is of the form $C \times J$, then $\mathcal{N}$ is of the form $\mathcal{M} \times J$, with $\mathcal{M} = \Ind_\pi(C)$, and the inclusion (4.10.3) is none other (up to isomorphism) than
\[
\underline{\operatorname{Hom}}(J, C) \longrightarrow \underline{\operatorname{Hom}}(J, \mathcal{M}) \ .
\]
Returning to the general case, we also see that:
\begin{enumerate}[label={}, leftmargin=3em]
\item[1°)] The second member of (4.10.3) is stable under $\varinjlim$ large relative to $\pi$, which are computed fiber by fiber.
\item[2°)] The first member of (4.10.3) consists of $\pi$-accessible objects of the second.
\end{enumerate}
Hence a \textit{fully} faithful functor
\[
\Ind_\pi(\uGamma(\mathcal{D}/J)) \hookrightarrow \uGamma(\mathcal{N}/J) \ . \tag*{(4.10.4)}
\]
\setcounter{footnote}{103}%
(\footnote{\texttt{[The numbers 4.10.1 to 4.10.4 are used twice.]}}).

\printedpage{92}
% head: A. Grothendieck : Dérivateurs, Ch. XVIII, Édition des Universités de Montpellier II et Paris VII
% foot: Version 17 décembre 2025

\textbf{Theorem 4.10.1.} \textit{Suppose that we have}
\[
\mathrm{card}\,\textup{red}\, J \ \leqslant \ \pi \ , \tag*{(4.10.1.1)}
\]
\textit{and that there exists a limit ordinal $\alpha$, such that:}
\begin{enumerate}[label={}, leftmargin=3em]
\item[\textit{1)}] $\card \alpha \leqslant \pi$.
\item[\textit{2)}] \textit{The direct limits of}
\msp{page 245}
\textit{type $I_\alpha = \{\beta \mid \beta < \alpha\}$ are representable in the categories $\mathcal{D}_j$ ($j \in \Ob J$), and for every arrow $\alpha : j \longrightarrow j'$ in $J$ \texttt{[too many $\alpha$]}, $\alpha^* : \mathcal{D}_j \longrightarrow \mathcal{D}_{j'}$ commutes with the said limits.}
\end{enumerate}
\textit{Then the canonical fully faithful functor (4.10.4) is an equivalence of categories.}

This will be contained in the slightly more precise result:

\textbf{Corollary 4.10.2.} \textit{The hypotheses being those of 4.10.1, with in addition}
\[
\card J \ \leqslant \ \pi \tag*{(4.10.2.1)}
\]
\textit{(which is equivalent to (4.10.1.1) if $J$ is reduced, a case to which we can of course reduce in 4.10.1). Let $\varphi \in \Ob \uGamma(\mathcal{N}/J)$, and let us consider the canonical} faithful \textit{functor:}
\[
\uGamma(\mathcal{D}/J)/\varphi \longrightarrow \prod_{j\in J} C_j/\varphi(j) \ . \tag*{(4.10.2.2)}
\]
\textit{Then the source category $\uGamma(\mathcal{D}/J)/\varphi$ is large relative to $\pi$, and the functor in question is cofinal.}

\msp{page 246}
From this, the theorem follows in an essentially sorital way, just as in the preceding special case 4.9.3 (cf. page 223) (\footnote{A bit short\,! The argument ought to be spelled out with care – there are two or three pages' worth of it\,!}).

To establish the corollary, since the target category of (4.10.2.2) is large relative to $\pi$, as a product of such categories (cf. 4.9.6.6, page 228), and the functor faithful, it suffices by virtue of 4.9.6.3, 2°, to verify conditions F 1 and F 3$_\pi$ (page 225) for the functor in question.

We are going to interpret the functor (4.10.2.2) in terms of a cofibered category
\[
\mathcal{D}_\varphi \longrightarrow J \ , \tag*{(4.10.5)}
\]
whose fibers are the $\mathcal{D}_j/\varphi(j)$. We note that if $\alpha : j \longrightarrow j'$ is an arrow of $J$, then we have
\[
\left\{
\begin{aligned}
&\alpha_* : C_j \ \longrightarrow \ C_{j'} \\
&\text{and} \ \ \alpha_*(\varphi(j)) \ \longrightarrow \ \varphi(j') \ ,
\end{aligned}
\right.
\]

\printedpage{93}
% head: A. Grothendieck : Dérivateurs, Ch. XVIII, Édition des Universités de Montpellier II et Paris VII
% foot: Version 17 décembre 2025

\msp{page 247}
hence a composite functor
\[
C_j/\varphi(j) \ \longrightarrow \ C_{j'}/\alpha^*(\varphi(j)) \ \longrightarrow \ C_{j'}/\varphi(j') \ ,
\]
with transitivity up to canonical isomorphism etc. We note that we have a canonical isomorphism
\[
\uGamma(\mathcal{D}_\varphi/J) \ \simeq \ \uGamma(\mathcal{D}/J)/\varphi \ ,
\]
so that the faithful canonical functor (4.10.2.2) is identified with the functor
\[
\uGamma(\mathcal{D}_\varphi/J) \longrightarrow \prod_{j\in\Ob J} (\mathcal{D}_\varphi)_j \ . \tag*{(4.10.6)}
\]

Corollary 4.10.2 will then be a consequence of the following more general result:

\textbf{Corollary 4.10.3.} \textit{Let $\mathcal{E} \longrightarrow J$ be a cofibrant functor, and $\pi$ an infinite cardinal. We assume:}
\begin{enumerate}[label={}, leftmargin=3em]
\item[\textit{a)}] $\card J \leqslant \pi$.
\item[\textit{b)}] \textit{The fiber categories $\mathcal{E}_j$ are equivalent to small categories, and large relative to $\pi$.}
\item[\textit{c)}] \textit{There exists a limit ordinal such that:}
\begin{enumerate}[label={}, leftmargin=3em]
\item[\textit{1°)}] $\card \alpha \leqslant \pi$.
\item[\textit{2°)}] \msp{page 248}\textit{The direct limits of type $I_\alpha$ are representable in the categories $\mathcal{E}_j$, and the cobase change functors $\lambda_* : \mathcal{E}_j \longrightarrow \mathcal{E}_{j'}$ commute with them. (This is the case e.g. if the} countable \textit{filtering direct limits are representable in the $\mathcal{E}_j$, and if the $\alpha_*$ commute with them\,; or if there exists a cardinal $\pi_0 \leqslant \pi$ and a cardinal $\pi'_0 > \pi_0$, such that the direct limits large relative to $\pi_0$ and of cardinality $\leqslant \pi'_0$ are representable in the $\mathcal{E}_j$, and the $\alpha_*$ commute with them.)}
\end{enumerate}
\end{enumerate}
\textit{Under these conditions, the faithful canonical functor}
\[
\uGamma(\mathcal{E}/J) \longrightarrow \prod_{j\in\Ob J} \mathcal{E}_j \tag*{(4.10.3.1)}
\]
\textit{is cofinal, and the source category is large relative to $\pi$.}

It is $\pm$ immediate that the hypotheses of 4.10.1 imply that those of 4.10.3 are satisfied for $\mathcal{E} = \mathcal{D}_\varphi \longrightarrow J$ (4.10.5). For
\msp{page 249}
this, we use the

% marginal vertical bar printed in the left margin beside the following statement
\textbf{Lemma 4.10.4.} \textit{Let $C$ be a category, and $\pi$ an infinite cardinal. Then the inclusion functor $C \longrightarrow \mathcal{M} \overset{\textup{def}}{=} \Ind_\pi C$ commutes with $\varinjlim$ of cardinality $\leqslant \pi$.}

Indeed, let $I$ be such that $\card I \leqslant \pi$, and $(x_i)_{i\in I}$ a direct system in $C$, having a limit $x$ there. On the other hand, let $X = \varinjlim\limits_{I'} \xi_{i'}$ be an object of $\Ind_\pi C$, with $I'$ large relative to $\pi$. We have

\end{document}
